\documentclass[conference,a4paper]{IEEEtran}

% ---------------------------------------------------------------
% Packages
% ---------------------------------------------------------------
\usepackage[T1]{fontenc}
\usepackage[utf8]{inputenc}
\usepackage{cite}
\usepackage{amsmath,amssymb,amsfonts}
\usepackage{graphicx}
\usepackage{textcomp}
\usepackage{xcolor}
\usepackage{hyperref}
\usepackage{booktabs}
\usepackage{array}
\usepackage{tabularx}

% ---------------------------------------------------------------
% Title & Authors  — PREENCHER PELO GRUPO
% ---------------------------------------------------------------
\title{Secure Software Lifecycle for a Portuguese Mobile Banking Application}

\author{
  \IEEEauthorblockN{Rafael Feliciano Soares}
  \IEEEauthorblockA{\textit{Mestrado em Cibersegurança}\\
    Universidade de Aveiro, Portugal\\
    rafael@ua.pt}
  \and
  \IEEEauthorblockN{Guilherme Almiro Viana\\ Fernandes Tavares}
  \IEEEauthorblockA{\textit{Mestrado em Cibersegurança}\\
    Universidade de Aveiro, Portugal\\
    gavftavares@ua.pt}
  \and
  \IEEEauthorblockN{Francisco Antunes dos Reis}
  \IEEEauthorblockA{\textit{Mestrado em Cibersegurança}\\
    Universidade de Aveiro, Portugal\\
    franciscoareis@ua.pt}
  \and
  \IEEEauthorblockN{Pedro Miguel Paulo Martins}
  \IEEEauthorblockA{\textit{Mestrado em Cibersegurança}\\
    Universidade de Aveiro, Portugal\\
    pedropmartins@ua.pt}
  \and
  \IEEEauthorblockN{Pedro Miguel \\dos Santos Marques}
  \IEEEauthorblockA{\textit{Mestrado em Cibersegurança}\\
    Universidade de Aveiro, Portugal\\
    pmsmarques@ua.pt}
}

\begin{document}
\maketitle

% ---------------------------------------------------------------
% ABSTRACT  (~10% do paper)
% ---------------------------------------------------------------
\begin{abstract}
% TODO: 1 parágrafo (~150 palavras).
% Diz: (1) o que é o sistema, (2) que abordagem de ciclo de vida seguro aplicaste,
% (3) qual a conclusão principal (ex.: os riscos mais críticos e como os mitigaste).
%
% Exemplo de arranque:
% "This paper applies a Secure Software Development Lifecycle (SSDL)
%  to a Portuguese mobile banking application offering transfers,
%  MB WAY integration and card controls. We map each lifecycle phase
%  to concrete security activities aligned with the NIST SSDF (SP~800-218)
%  and OWASP SAMM v2 frameworks, derive six measurable non-functional
%  security requirements (NFRs), and identify the three riskiest design
%  violations with their mitigations. Results show that ..."
\end{abstract}

\begin{IEEEkeywords}
secure software lifecycle, SSDL, mobile banking, NIST SSDF, OWASP SAMM, NFR, threat modelling
\end{IEEEkeywords}

% ---------------------------------------------------------------
% 1. INTRODUCTION  (~30% do paper)
% ---------------------------------------------------------------
\section{Introduction}
\label{sec:intro}

% TODO: Apresenta o sistema — responde a:
%   - O que é? (mobile app banco português, transferências, MB WAY, controlos cartão)
%   - Quem são os utilizadores e interfaces? (~2-3 interfaces)
%   - Que funções cobre? (5-8 funções, foco em cibersegurança)
%   - Porquê é crítico? (dados financeiros, regulação PSD2/DORA, fraude)
%   - Estrutura do paper (1 parágrafo no fim da intro)

The system under analysis is a mobile banking application for a Portuguese retail bank,
providing customers with core financial services including fund transfers (SEPA and
instant), MB~WAY integration, card controls (freeze, spending limits, PIN change), and
account statement access. The application interacts with three main external interfaces:
the bank's core banking system via an internal REST~API, the MB~WAY interoperability
network, and the Portuguese national payment scheme (Multibanco/SIBS).

\subsection{Security Context}

% TODO: regulação aplicável (PSD2, DORA, GDPR) e ameaças típicas (phishing, MitM, credential stuffing)
% ~2 parágrafos

Mobile banking applications are subject to strict regulatory requirements in Portugal and
the European Union, including the Payment Services Directive~2 (PSD2)~\cite{psd2},
the Digital Operational Resilience Act (DORA)~\cite{dora}, and the General Data
Protection Regulation (GDPR). PSD2 mandates Strong Customer Authentication (SCA) for
all payment transactions, while DORA requires financial entities to demonstrate
operational resilience and incident response capability. Given the high value of
financial assets and sensitive personal data handled, such applications face a broad
threat landscape including credential stuffing, phishing, man-in-the-middle attacks,
and API abuse.

\subsection{Scope and Approach}

% TODO: diz que vais aplicar SSDL fase a fase, usando NIST SSDF como referência principal

This paper applies a Secure Software Development Lifecycle (SSDL) to the mobile
banking application, mapping security activities to each development phase following
the NIST Secure Software Development Framework (SSDF, SP~800-218)~\cite{nist_ssdf}
and drawing on OWASP SAMM~v2~\cite{owasp_samm} for maturity assessment. The remainder
of the paper is organised as follows: Section~\ref{sec:lifecycle} describes the
security activities per lifecycle phase; Section~\ref{sec:nfr} presents six measurable
non-functional security requirements; Section~\ref{sec:principles} evaluates the
system against ten secure design principles; Section~\ref{sec:conclusion} concludes.

% ---------------------------------------------------------------
% 2. SECURE SOFTWARE LIFECYCLE  (~50% do paper — PARTE PRINCIPAL)
% ---------------------------------------------------------------
\section{Secure Software Lifecycle}
\label{sec:lifecycle}

% Guia: para cada fase, descreve (1) o que se faz em geral e (2) como se aplica
% especificamente à app de banking. ~1-2 parágrafos + lista de actividades por fase.

\subsection{Training}

% TODO: que formação a equipa de dev precisa? Ex.: OWASP Mobile Top 10, secure coding
% em Swift/Kotlin, PSD2/SCA, tratamento de dados financeiros (GDPR)
% Output esperado: training records (evidência para auditoria)

Before development begins, all team members must complete role-based secure development
training. For a mobile banking application the relevant topics include: OWASP Mobile
Security Testing Guide (MSTG)~\cite{owasp_mstg}, secure coding in Android (Kotlin) and
iOS (Swift), cryptographic key management, PSD2/SCA implementation requirements, and
GDPR data handling obligations.

\subsection{Requirements}

% TODO: security requirements específicos do banking, abuse cases, compliance mapping
% Output esperado: security requirements catalogue
% Dica: liga a Section 3 (NFRs) aqui — "the measurable NFRs are defined in Section III"

Security requirements are derived from regulatory obligations (PSD2, DORA, GDPR),
known threat patterns (OWASP Mobile Top~10), and stakeholder abuse cases. Key
requirements include: enforcement of Strong Customer Authentication for all payment
operations, end-to-end encryption of all financial data in transit and at rest, session
timeout and re-authentication policies, and tamper-evident audit logging of all
privileged operations. Measurable non-functional requirements (NFRs) derived from these
are presented in Section~\ref{sec:nfr}.

\subsection{Design}

% TODO: threat modelling (STRIDE sobre o contexto diagram), decisões de arquitectura
% (ex: certificate pinning, token-based auth, offline vs online), crypto choices
% Output esperado: threat model, design review minutes

The design phase centres on threat modelling using the STRIDE methodology applied to the
system's context diagram. The main components — mobile client, API gateway, core banking
backend, and MB~WAY integration — are analysed for Spoofing, Tampering, Repudiation,
Information Disclosure, Denial of Service, and Elevation of Privilege threats.

% TODO: descreve 3-4 decisões de arquitectura concretas e porquê (ex: certificate
% pinning para evitar MitM, JWT com curta expiração, separação de ambientes prod/dev)

Key architectural decisions include: TLS~1.3 with certificate pinning on the mobile
client to prevent man-in-the-middle attacks; short-lived JWT access tokens (15~min
expiry) with refresh token rotation; strict separation of production and development
environments (NIST SSDF PO.5); and a zero-trust API gateway enforcing per-endpoint
authorisation.

\subsection{Implementation}

% TODO: secure coding standard adoptado, SAST tool usado, dependency scanning,
% secret scanning no CI/CD pipeline
% Output esperado: clean SAST gate, SBOM

During implementation the team follows the OWASP Mobile Application Security
Verification Standard (MASVS)~\cite{owasp_masvs} as the secure coding standard.
Static Application Security Testing (SAST) is integrated into the CI/CD pipeline,
blocking merges when high-severity findings are detected. Software Composition Analysis
(SCA) scans all third-party dependencies for known vulnerabilities (CVEs), and secret
scanning prevents credentials from being committed to the repository. A Software Bill
of Materials (SBOM) is generated at each release.

\subsection{Verification}

% TODO: DAST, fuzzing da API, pentest da app móvel, code review focalizado em auth
% Output esperado: test reports, residual-risk list

Verification combines automated and manual techniques. Dynamic Application Security
Testing (DAST) is performed against the staging API environment. The mobile application
undergoes penetration testing following the OWASP MSTG checklist, covering reverse
engineering protections, insecure data storage, and authentication bypass. Fuzz testing
targets the payment API endpoints. All findings are recorded in a residual-risk list
with severity, owner, and remediation deadline.

\subsection{Release}

% TODO: final security review, go/no-go criteria, signing of artefacts, incident response plan
% Output esperado: go/no-go record, response runbook

Before production release, a final security review validates that all critical and
high-severity findings are resolved and the residual-risk list has been accepted by the
security officer. Release artefacts (APK/IPA) are signed with the organisation's
code-signing certificate, and a SHA-256 hash is published for integrity verification.
An incident response runbook specific to the banking application is prepared and
reviewed with the PSIRT team.

\subsection{Operations}

% TODO: monitorização runtime (ex: anomaly detection em transacções), vulnerability
% disclosure, patching cadence, EoL plan
% Output esperado: advisories, patch cadence, PSIRT

In production, the application is monitored for anomalous transaction patterns and
authentication failures using a Security Information and Event Management (SIEM) system.
A vulnerability disclosure policy is published, and a patch cadence of 30~days for
critical CVEs and 90~days for medium severity is defined. An end-of-life plan ensures
that deprecated API versions are retired with adequate user notice.

% ---------------------------------------------------------------
% 3. MEASURABLE NFRs  (~parte da secção de lifecycle)
% ---------------------------------------------------------------
\section{Measurable Non-Functional Requirements}
\label{sec:nfr}

% TODO: 3 atributos de qualidade escolhidos + justificação, depois 2 NFRs por atributo.
% Formato de um bom NFR: condição + comportamento mensurável + threshold + método de verificação
% Exemplo: "When a user enters an incorrect PIN more than 3 times within 60~s,
%           the system shall lock the account for 15~min and generate an audit event
%           (Security -- Accountability; verified by functional test T-AUTH-02)."

Based on the system's risk profile, three quality attributes were selected as most
critical: \textbf{Security} (confidentiality and accountability sub-characteristics),
\textbf{Availability}, and \textbf{Integrity}.

\subsection{Security}

\begin{itemize}
  \item \textbf{NFR-SEC-01 (Accountability):} When a user performs any fund transfer,
    the system shall record a tamper-evident audit event containing user~ID, timestamp,
    source and destination accounts, and amount, within 500~ms of the operation.
    \textit{Verification: automated audit-log test suite (T-SEC-01).}

  \item \textbf{NFR-SEC-02 (Confidentiality):} A user without the ``card-manager'' role
    shall not be able to retrieve another user's card~PIN or card number in plain text
    via any API endpoint.
    \textit{Verification: negative access-control test (T-SEC-02).}
\end{itemize}

\subsection{Availability}

\begin{itemize}
  \item \textbf{NFR-AVL-01:} The payment API shall maintain a monthly availability of
    $\geq$99.9\% (excluding planned maintenance windows of $\leq$4~h/month), degrading
    gracefully to read-only mode when the core banking backend is unreachable.
    \textit{Verification: availability monitoring dashboard + chaos engineering test (T-AVL-01).}

  \item \textbf{NFR-AVL-02:} After a watchdog-triggered service restart, the API
    gateway shall restore full service within $\leq$60~s.
    \textit{Verification: restart recovery test in staging (T-AVL-02).}
\end{itemize}

\subsection{Integrity}

\begin{itemize}
  \item \textbf{NFR-INT-01:} Any API response carrying a transaction record shall
    include an HMAC-SHA256 signature; the mobile client shall reject responses whose
    signature does not verify against the published public key.
    \textit{Verification: tamper-injection test (T-INT-01).}

  \item \textbf{NFR-INT-02:} Release artefacts (APK/IPA) shall be code-signed, and
    the mobile OS shall reject unsigned or improperly signed builds at install time.
    \textit{Verification: install test with a tampered binary (T-INT-02).}
\end{itemize}

% ---------------------------------------------------------------
% 4. DESIGN PRINCIPLES CHECKLIST  (~10% do paper)
% ---------------------------------------------------------------
\section{Secure Design Principles Evaluation}
\label{sec:principles}

% TODO: preencher a tabela com os 10 princípios da S01.
% Para cada um: honra (H) ou viola (V), e uma linha de justificação.
% Os 3 piores violations têm de ter mitigação (última coluna).
% Ajusta os princípios se a S01 tiver nomes ligeiramente diferentes.

Table~\ref{tab:principles} evaluates the banking application against the ten secure
design principles introduced in RS-S01.

\begin{table*}[t]
\caption{Secure Design Principles Checklist}
\label{tab:principles}
\begin{tabularx}{\textwidth}{lXcX}
\toprule
\textbf{\#} & \textbf{Principle} & \textbf{H/V} & \textbf{Justification / Mitigation} \\
\midrule
1 & Least Privilege & H &
  A regular user has no access to user management and cannot freeze or delete other accounts.
  An admin cannot initiate transfers or manage a user's account directly. \\
2 & Fail-safe Defaults & H &
  Failed login denies access. A failed transfer is automatically cancelled and never processed. \\
3 & Economy of Mechanism & H &
  A centralised permission manager serves both the mobile app and desktop interfaces,
  providing a single authoritative access-control point. \\
4 & Complete Mediation & \textbf{V} &
  Login credentials are cached; the user is not re-authenticated on every session entry.
  \textit{Mitigation: introduce a 4-digit PIN for fast re-authentication; enforce session
  timeout after 15~min of inactivity; require MFA only for high-value operations such
  as large transfers.} \\
5 & Open Design & H &
  Even knowing the bank's internal security architecture, an attacker cannot access a
  user account without the correct password/PIN. Security relies on key secrecy, not
  obscurity. \\
6 & Separation of Privilege & H &
  Large transfers and account creation require MFA. An admin freezing an account must
  obtain confirmation from a supervisor or a second admin. \\
7 & Least Common Mechanism & \textbf{V} &
  All data is stored on a single global server instead of region-isolated servers,
  creating a single point of failure and shared-risk exposure.
  \textit{Mitigation: partition critical-data components; isolate internal services;
  deploy independent regional storage for sensitive data.} \\
8 & Psychological Acceptability & \textbf{V} &
  Users reject frequent MFA prompts, mandatory login on every app launch, and
  rapidly expiring session tokens, which they consider overly intrusive.
  \textit{Mitigation: 4-digit PIN instead of a complex password; reasonable session
  duration with silent renewal; step-up MFA only when genuinely necessary.} \\
9 & Defence in Depth & H &
  Multiple independent layers: password, MFA via email/SMS, session tokens with
  expiry, and a bank-issued verification matrix for account creation or confirmation. \\
10 & Fail Securely & H &
  An error during a loan request displays only a generic message (``action could not
  be completed''); no technical details are exposed that could be exploited to bypass
  interest calculations or other business rules. \\
\bottomrule
\end{tabularx}
\end{table*}

% ---------------------------------------------------------------
% 5. CONCLUSION
% ---------------------------------------------------------------
\section{Conclusion}
\label{sec:conclusion}

% TODO: ~1-2 parágrafos.
% Resume: o que fizeste, os riscos mais críticos encontrados, e os próximos passos
% (Exercise 2 vai usar os NFRs daqui).

This paper applied a Secure Software Development Lifecycle to a Portuguese mobile
banking application, mapping concrete security activities to each phase from training
to operations, drawing on the NIST~SSDF and OWASP~SAMM frameworks. Three quality
attributes -- Security, Availability, and Integrity -- were selected as the most
critical for this domain, yielding six measurable NFRs that will seed the requirements
catalogue in the next exercise. The three most significant design violations identified are: incomplete mediation
(login credentials cached across sessions), least common mechanism (single global
server with no regional isolation), and psychological acceptability (users rejecting
frequent MFA and short session timeouts). Mitigations for each have been defined at
the design and requirements level and will seed the requirements catalogue in Exercise~2.

% ---------------------------------------------------------------
% REFERENCES
% ---------------------------------------------------------------
\begin{thebibliography}{00}

\bibitem{psd2}
European Parliament, ``Directive (EU) 2015/2366 on Payment Services (PSD2),''
\textit{Official Journal of the European Union}, Dec. 2015.

\bibitem{dora}
European Parliament, ``Regulation (EU) 2022/2554 -- Digital Operational Resilience Act
(DORA),'' \textit{Official Journal of the European Union}, Dec. 2022.

\bibitem{nist_ssdf}
M. Souppaya, K. Scarfone, and D. Dodson, ``Secure Software Development Framework
(SSDF) Version~1.1,'' NIST Special Publication~800-218, Feb. 2022.

\bibitem{owasp_samm}
OWASP Foundation, ``Software Assurance Maturity Model (SAMM) v2.0,'' 2020.
[Online]. Available: \url{https://owaspsamm.org}

\bibitem{owasp_mstg}
OWASP Foundation, ``Mobile Security Testing Guide (MSTG),'' 2023.
[Online]. Available: \url{https://mas.owasp.org/MASTG/}

\bibitem{owasp_masvs}
OWASP Foundation, ``Mobile Application Security Verification Standard (MASVS) v2,''
2023. [Online]. Available: \url{https://mas.owasp.org/MASVS/}

\end{thebibliography}

\end{document}
